% ASPLOS 2027 Paper Template
% ACM International Conference on Architectural Support for
% Programming Languages and Operating Systems
%
% Format: ACM SIGPLAN, ≤12 pages (excluding references), 10pt, two-column
% Uses acmart.cls with sigplan option
%
% Official CFP: https://www.asplos-conference.org/asplos2026/call-for-papers-asplos27/
% ACM Template: https://www.acm.org/publications/proceedings-template
%
% IMPORTANT NOTES:
% - RAPID REVIEW ROUND: Reviewers read ONLY first 2 pages!
%   → First 2 pages MUST be self-contained
%   → Clearly state problem, approach, contribution in first 2 pages
% - Must advance Architecture, PL, and/or OS research
% - Two cycles: April 2026 and September 2026
% - Max 4 submissions per author per cycle
% - Major Revision decision available
% - Double-blind review
%
% RAPID REVIEW TIPS:
% - Page 1: Problem motivation + why it matters to arch/PL/OS
% - Page 2: Approach overview + key results preview + contribution list
% - Do NOT rely on content after page 2 for rapid review

\documentclass[sigplan,10pt]{acmart}

% Remove copyright/permission footer for submission
\renewcommand\footnotetextcopyrightpermission[1]{}
\settopmatter{printfolios=true}

% Remove ACM reference format for submission
\setcopyright{none}
\renewcommand\acmConference[4]{}
\acmDOI{}
\acmISBN{}

% Recommended packages
\usepackage{booktabs}
\usepackage{xspace}
\usepackage{subcaption}

% Custom commands
\newcommand{\system}{SystemName\xspace}  % Replace with your anonymized system name
\newcommand{\eg}{e.g.,\xspace}
\newcommand{\ie}{i.e.,\xspace}
\newcommand{\etal}{\textit{et al.}\xspace}

\begin{document}

\title{Your Paper Title Here}

% Anonymized for submission
\author{Paper \#XXX}
\affiliation{}

% Camera-ready:
% \author{Author One}
% \affiliation{%
%   \institution{University/Company}
%   \city{City}
%   \country{Country}}
% \email{email@example.com}

\begin{abstract}
% Write your abstract here.
% Remember: abstract + first 2 pages are critical for rapid review.
Your abstract goes here.
\end{abstract}

\maketitle
\pagestyle{plain}

\section{Introduction}
\label{sec:intro}

% CRITICAL FOR RAPID REVIEW (Pages 1-2):
% Reviewers will read ONLY the first 2 pages in the rapid review round.
% These 2 pages must:
% 1. Clearly state the problem and why it matters
% 2. Explain how this advances Architecture, PL, or OS research
%    (NOT just using arch/PL/OS to advance another domain)
% 3. Outline your approach and key contributions
% 4. Preview your main results
%
% If reviewers cannot determine your contribution to arch/PL/OS
% from the first 2 pages, your paper may be rejected in rapid review.

Your introduction goes here. Make the first 2 pages self-contained.
Clearly articulate how this work advances computer architecture,
programming languages, or operating systems research.

\section{Background and Motivation}
\label{sec:background}

Provide necessary background.

\section{Design}
\label{sec:design}

Describe the design of your system/approach.

\section{Implementation}
\label{sec:implementation}

Describe implementation details.

\section{Evaluation}
\label{sec:evaluation}

Present your evaluation results.

\subsection{Experimental Setup}

Describe your experimental setup.

\subsection{Results}

Present and discuss your results.

\section{Related Work}
\label{sec:related}

Discuss related work.

\section{Conclusion}
\label{sec:conclusion}

Summarize your contributions.

\bibliographystyle{ACM-Reference-Format}
\bibliography{references}

\end{document}
